\section{Topological spaces}

\subsection{Definitions and examples}

\begin{de}
A topological space is a pair $(X,\mathcal{T})$, where $X$ is a set and
$\mathcal{T}$ is a collection of subsets of $X$ satisfying:

(a) $\varnothing,X\in\mathcal{T}$;

(b) if $\{G_i:i\in I\}\subseteq\mathcal{T}$, then
$\bigcup_{i\in I}G_i\in\mathcal{T}$;

(c) if $G_1,\ldots,G_n\in\mathcal{T}$, then
$\bigcap_{k=1}^nG_k\in\mathcal{T}$.

The collection $\mathcal{T}$ is called a topology on $X$, and its members are
called open sets.

For any set $X$, the collection
\[
\mathcal{T}=\{G\subseteq X:X\setminus G\text{ is finite}\}\cup\{\varnothing\}
\]
is the \emph{cofinite topology}. If $X$ is infinite, this topology is not
Hausdorff.

If $(X,\mathcal{T})$ is a topological space and $Y\subseteq X$, then
\[
\mathcal{T}_Y=\{Y\cap G:G\in\mathcal{T}\}
\]
is a topology on $Y$, called the subspace topology or relative topology. For
a metric space, this agrees with the topology induced by the restricted
metric.

A topological space $X$ is Hausdorff if, whenever $x\neq y$, there are
disjoint open sets $U$ and $V$ such that $x\in U$ and $y\in V$.
\end{de}

\textbf{Agreement.} Unless explicitly stated otherwise, all topological
spaces below are assumed to be Hausdorff.

A subset $F$ of a topological space $X$ is closed if $X\setminus F$ is open.

\begin{thm}
Let $X$ be a topological space, and let $\mathcal{F}$ be the collection of
closed subsets of $X$.

(a) $\varnothing,X\in\mathcal{F}$.

(b) An arbitrary intersection of members of $\mathcal{F}$ belongs to
$\mathcal{F}$.

(c) A finite union of members of $\mathcal{F}$ belongs to $\mathcal{F}$.

(d) If $X$ is Hausdorff, then every singleton $\{x\}$ is closed.
\end{thm}

\begin{proof}
Parts (a)--(c) follow from De Morgan's laws. For (d), fix $x\in X$. For every
$y\neq x$, choose an open set $U_y$ containing $y$ and disjoint from some open
neighborhood of $x$. Then
$X\setminus\{x\}=\bigcup_{y\neq x}U_y$ is open.
\end{proof}

\begin{de}
Let $A\subseteq X$. Its interior, closure, and boundary are
\[
\operatorname{int}A=\bigcup\{G:G\text{ is open and }G\subseteq A\},
\]
\[
\operatorname{cl}A=\bigcap\{F:F\text{ is closed and }A\subseteq F\},
\]
and
\[
\partial A=\operatorname{cl}A\cap\operatorname{cl}(X\setminus A).
\]
\end{de}

\begin{thm}
Let $A\subseteq X$.

(a) $x\in\operatorname{int}A$ if and only if there is an open set $G$ such
that $x\in G\subseteq A$.

(b) $x\in\operatorname{cl}A$ if and only if every open set $G$ containing
$x$ meets $A$.

(c) $\operatorname{int}A$ is the largest open subset of $A$.

(d) $\operatorname{cl}A$ is the smallest closed set containing $A$.
\end{thm}

\begin{proof}
All four statements follow directly from the definitions.
\end{proof}

\begin{thm}
Let $A\subseteq X$.

(a) $A$ is closed if and only if $A=\operatorname{cl}A$.

(b) $A$ is open if and only if $A=\operatorname{int}A$.

(c) If $A_1,\ldots,A_n\subseteq X$, then
\[
\operatorname{cl}\left(\bigcup_{k=1}^nA_k\right)
=\bigcup_{k=1}^n\operatorname{cl}A_k.
\]
\end{thm}

\begin{proof}
Parts (a) and (b) follow from the extremal properties in the preceding
theorem. For (c), the right-hand side is a closed set containing the union,
which gives one inclusion. Conversely, each $A_k$ is contained in the union,
so each $\operatorname{cl}A_k$ is contained in the closure of the union.
\end{proof}

\begin{de}
A subset $E$ of a topological space is dense if
$\operatorname{cl}E=X$. A topological space is separable if it has a
countable dense subset.
\end{de}

\begin{thm}
A set $E$ is dense in $X$ if and only if every nonempty open subset of $X$
meets $E$. Equivalently, every open neighborhood of every point meets $E$.
\end{thm}

\begin{de}
A point $x\in X$ is a limit point of $A\subseteq X$ if every open
neighborhood $G$ of $x$ satisfies
\[
(G\setminus\{x\})\cap A\neq\varnothing.
\]
\end{de}

\begin{thm}
Let $A\subseteq X$.

(a) The set $A$ is closed if and only if it contains all its limit points.

(b)
\[
\operatorname{cl}A=A\cup\{x:x\text{ is a limit point of }A\}.
\]
\end{thm}

\begin{proof}
If $x\in\operatorname{cl}A\setminus A$, then every neighborhood of $x$ meets
$A$, and any such point of $A$ is different from $x$; hence $x$ is a limit
point. Conversely, every limit point belongs to $\operatorname{cl}A$ because
every neighborhood of it meets $A$. This proves (b), and (a) follows.
\end{proof}

\begin{thm}
Let $Y\subseteq X$ have the subspace topology.

(a) A subset $A\subseteq Y$ is closed in $Y$ if and only if
$A=F\cap Y$ for some closed set $F\subseteq X$.

(b) If $A\subseteq Y$, then
\[
\operatorname{cl}_Y A=Y\cap\operatorname{cl}_X A.
\]

(c) If $A\subseteq Y$, then
\[
Y\cap\operatorname{int}_X A\subseteq\operatorname{int}_Y A.
\]
The inclusion can be strict.
\end{thm}

\begin{proof}
(a) The set $A$ is closed in $Y$ if and only if $Y\setminus A=Y\cap G$ for
some open $G\subseteq X$. This is equivalent to
$A=Y\cap(X\setminus G)$.

(b) The set $Y\cap\operatorname{cl}_X A$ is closed in $Y$ and contains $A$.
If $C$ is any closed subset of $Y$ containing $A$, write $C=Y\cap F$ with
$F$ closed in $X$. Then $\operatorname{cl}_X A\subseteq F$, so
$Y\cap\operatorname{cl}_X A\subseteq C$.

(c) If $x\in Y\cap\operatorname{int}_X A$, choose an open set
$G\subseteq X$ with $x\in G\subseteq A$. Then $G\cap Y$ is an open
neighborhood of $x$ in $Y$ contained in $A$.
\end{proof}

\subsection{Bases and subbases}

\begin{de}
A collection $\mathcal{B}$ of subsets of $X$ is a base for a topology
$\mathcal{T}$ if every open set is a union of members of $\mathcal{B}$.
\end{de}

\begin{thm}
A collection $\mathcal{B}$ of subsets of $X$ is a base for a topology on $X$
if and only if:

(a) $X=\bigcup_{B\in\mathcal{B}}B$;

(b) whenever $B_1,B_2\in\mathcal{B}$ and $x\in B_1\cap B_2$, there is a
$B\in\mathcal{B}$ such that
$x\in B\subseteq B_1\cap B_2$.

The topology generated by $\mathcal{B}$ is Hausdorff if and only if, in
addition,

(c) for distinct $x,y\in X$, there are disjoint $A,B\in\mathcal{B}$ with
$x\in A$ and $y\in B$.
\end{thm}

\begin{proof}
If $\mathcal{B}$ is a base, (a) and (b) follow from the fact that $X$ and
$B_1\cap B_2$ are open. Conversely, let $\mathcal{T}$ be the collection of
all unions of members of $\mathcal{B}$. Conditions (a) and (b) give the empty
set, the whole space, arbitrary unions, and finite intersections required of
a topology. Condition (c) is exactly the Hausdorff separation condition
expressed using basic neighborhoods.
\end{proof}

\begin{cor}
Let $(X,d)$ be separable, let $\{a_1,a_2,\ldots\}$ be dense in $X$, and let
$\{r_1,r_2,\ldots\}$ enumerate the positive rational numbers. Then
\[
\mathcal{B}=\{B(a_n;r_m):n,m\geq1\}
\]
is a countable base for the metric topology.
\end{cor}

\begin{proof}
Let $G$ be open and $x\in G$. Choose $\varepsilon>0$ with
$B(x;\varepsilon)\subseteq G$. Pick $a_n$ with
$d(x,a_n)<\varepsilon/3$, and choose a positive rational $r_m$ satisfying
\[
d(x,a_n)<r_m<\varepsilon-d(x,a_n).
\]
Then $x\in B(a_n;r_m)\subseteq G$.
\end{proof}

\begin{cor}
Let $\mathcal{S}$ be any collection of subsets of $X$. The collection
$\mathcal{B}$ of all finite intersections of members of $\mathcal{S}$,
including the empty intersection $X$, is a base for a topology on $X$.
This topology is Hausdorff if $\mathcal{S}$ separates points in the following
sense: for every $x\neq y$, there are disjoint $S,T\in\mathcal{S}$ with
$x\in S$ and $y\in T$.
\end{cor}

\begin{proof}
Finite intersections cover $X$ because the empty intersection is $X$, and
the intersection of two finite intersections is again a finite intersection.
The final assertion follows from part (c) of the basis theorem.
\end{proof}

\begin{de}
A collection $\mathcal{S}$ of subsets of $X$ is called a subbase if the finite
intersections of its members form a base. The resulting topology is the
topology generated by $\mathcal{S}$.
\end{de}

\begin{ex}
A partially ordered set is a pair $(X,\leq)$ in which $\leq$ is reflexive,
transitive, and antisymmetric. It is linearly ordered if any two points are
comparable.

For $a\in X$, put
\[
(-\infty,a)=\{x\in X:x<a\},\qquad
(a,\infty)=\{x\in X:a<x\}.
\]
Let $\mathcal{S}$ be the collection of all such rays.

(a) Show that $\mathcal{S}$ generates a topology on every partially ordered
set. This topology need not be Hausdorff.

(b) Show that it is Hausdorff when the order is linear.

(c) In the linearly ordered case, show that a base consists of the open
intervals $(a,b)$ together with the rays $(-\infty,b)$ and $(a,\infty)$ that
are needed at endpoints.
\end{ex}

\begin{proof}
Part (a) follows from the subbase construction. For (b), let $x<y$. If there
is a point $z$ with $x<z<y$, then $(-\infty,z)$ and $(z,\infty)$ separate
$x$ and $y$. If no such $z$ exists, then $(-\infty,y)$ and $(x,\infty)$ do.
Part (c) follows by listing the possible nonempty finite intersections of
left and right rays.
\end{proof}

\subsection{Continuous functions}

\begin{de}
Let $X$ and $W$ be topological spaces. A function $f:X\to W$ is continuous
at $x\in X$ if, for every neighborhood $U$ of $f(x)$, there is a neighborhood
$G$ of $x$ such that $f(G)\subseteq U$. It is continuous if it is continuous
at every point.
\end{de}

\begin{thm}
For a function $f:X\to W$, the following are equivalent:

(a) $f$ is continuous;

(b) $f^{-1}(U)$ is open in $X$ for every open $U\subseteq W$;

(c) $f^{-1}(C)$ is closed in $X$ for every closed $C\subseteq W$;

(d) if $\mathcal{B}$ is a base for $W$, then $f^{-1}(B)$ is open for every
$B\in\mathcal{B}$;

(e) if $\mathcal{S}$ is a subbase for $W$, then $f^{-1}(S)$ is open for every
$S\in\mathcal{S}$.
\end{thm}

\begin{proof}
The equivalence of (a) and (b) follows from the neighborhood definition. The
equivalence of (b) and (c) follows because inverse images commute with
complements. Clearly (b) implies (d), and (d) implies (e) because subbasic
sets are open. Finally, if (e) holds, inverse images of finite intersections
of subbasic sets are open, and inverse images commute with arbitrary unions;
therefore the inverse image of every open set is open.
\end{proof}

\begin{thm}
The composition of continuous functions is continuous.
\end{thm}

\begin{proof}
If $f:X\to Y$ and $g:Y\to W$ are continuous and $U\subseteq W$ is open,
then
\[
(g\circ f)^{-1}(U)=f^{-1}(g^{-1}(U))
\]
is open.
\end{proof}

\begin{thm}[Pasting Lemma]
Let $X=A\cup B$, and let $g:A\to W$ and $h:B\to W$ be continuous functions
that agree on $A\cap B$. Define $f:X\to W$ by $f|_A=g$ and $f|_B=h$. If
$A$ and $B$ are both open, or if they are both closed, then $f$ is
continuous.
\end{thm}

\begin{proof}
If $A$ and $B$ are open and $U\subseteq W$ is open, then $g^{-1}(U)$ is open
in $A$ and hence in $X$; similarly for $h^{-1}(U)$. Thus
$f^{-1}(U)=g^{-1}(U)\cup h^{-1}(U)$ is open.

If $A$ and $B$ are closed, apply the same argument to inverse images of
closed sets: a set closed in $A$ is closed in $X$ because $A$ is closed, and
similarly for $B$.
\end{proof}

\begin{thm}
Let $X_k$ be a topological space for $1\leq k\leq n$, and put
$X=X_1\times\cdots\times X_n$. The collection
\[
\mathcal{B}=\{G_1\times\cdots\times G_n:G_k\text{ is open in }X_k\}
\]
is a base for a topology on $X$. This topology is the smallest topology for
which every coordinate projection $\pi_k:X\to X_k$ is continuous.
\end{thm}

\begin{proof}
The basis conditions are immediate. For an open set $U\subseteq X_k$,
\[
\pi_k^{-1}(U)=X_1\times\cdots\times X_{k-1}\times U\times
X_{k+1}\times\cdots\times X_n
\]
is basic open, so every projection is continuous. Conversely, if a topology
on $X$ makes all projections continuous, then
\[
G_1\times\cdots\times G_n=\bigcap_{k=1}^n\pi_k^{-1}(G_k)
\]
is open. Thus it contains the product topology.
\end{proof}

\begin{de}
The topology in the preceding theorem is the product topology. The maps
$\pi_k$ are the coordinate projections.
\end{de}

\begin{de}
A function $f:X\to Z$ is an open map if $f(G)$ is open in $Z$ whenever $G$
is open in $X$.
\end{de}

\begin{thm}
Every coordinate projection from a finite product with the product topology
is an open map.
\end{thm}

\begin{proof}
If the product is empty, the assertion is immediate. Otherwise, every factor
is nonempty, and the image under $\pi_k$ of a basic open set
$G_1\times\cdots\times G_n$ is either empty or $G_k$. Since every open set
is a union of basic open sets and projections commute with unions,
$\pi_k$ is open.
\end{proof}

\begin{thm}
Let $Y$ be a topological space and let
$f:Y\to X_1\times\cdots\times X_n$. Then $f$ is continuous if and only if
each coordinate function $\pi_k\circ f$ is continuous.
\end{thm}

\begin{proof}
One direction follows by composition. Conversely, for a basic open set
$G_1\times\cdots\times G_n$,
\[
f^{-1}(G_1\times\cdots\times G_n)
=\bigcap_{k=1}^n(\pi_k\circ f)^{-1}(G_k),
\]
which is open.
\end{proof}

\begin{thm}
If $f,g:X\to\mathbb{R}$ are continuous, then $f+g$ and $fg$ are continuous.
\end{thm}

\begin{proof}
The map $x\mapsto(f(x),g(x))$ is continuous by the product criterion. Compose
it with the continuous maps $(a,b)\mapsto a+b$ and $(a,b)\mapsto ab$.
\end{proof}

\begin{thm}
If $f,g:X\to\mathbb{R}$ are continuous, then so are
$f\vee g$, $f\wedge g$, and $|f|$.
\end{thm}

\begin{proof}
The absolute-value function is continuous, and
\[
f\vee g=\frac{f+g+|f-g|}{2},\qquad
f\wedge g=\frac{f+g-|f-g|}{2}.
\]
\end{proof}

\begin{de}
A homeomorphism between topological spaces $X$ and $Z$ is a bijection
$f:X\to Z$ such that both $f$ and $f^{-1}$ are continuous.
\end{de}

\subsection{Compactness and connectedness}

\begin{de}
A subset $K$ of a topological space $X$ is compact if every open cover of
$K$ has a finite subcover.
\end{de}

\begin{thm}
Let $X$ be Hausdorff and let $K\subseteq X$.

(a) If $K$ is compact, then $K$ is closed.

(b) If $K$ is compact and $F\subseteq K$ is closed in $K$, then $F$ is
compact.

(c) The continuous image of a compact set is compact.
\end{thm}

\begin{proof}
(a) Fix $x\notin K$. For each $z\in K$, choose disjoint open sets $U_z$ and
$V_z$ with $x\in U_z$ and $z\in V_z$. Finitely many $V_{z_1},\ldots,V_{z_n}$
cover $K$. Then $U=\bigcap_{k=1}^nU_{z_k}$ is an open neighborhood of $x$
disjoint from $K$. Hence $X\setminus K$ is open.

(b) By part (a), $K$ is closed in $X$. Since $F$ is closed in $K$, it is
closed in $X$. Add $X\setminus F$ to any open cover of $F$ to obtain an open
cover of $K$.

(c) Pull an open cover of $f(K)$ back under $f$, choose a finite subcover of
$K$, and push it forward.
\end{proof}

\begin{cor}
If $X$ is nonempty and compact and $f:X\to\mathbb{R}$ is continuous,
then $f$ attains its minimum and maximum.
\end{cor}

\begin{thm}
A subset $K$ of a topological space is compact if and only if every collection
of closed subsets of $K$ having the finite intersection property has
nonempty intersection.
\end{thm}

\begin{proof}
Take complements in $K$. A family of closed sets has the finite intersection
property exactly when the corresponding family of open complements has no
finite subcover; its total intersection is empty exactly when those
complements cover $K$.
\end{proof}

\begin{thm}
If $\mathcal{B}$ is a base for $X$, then $K\subseteq X$ is compact if and only
if every cover of $K$ by members of $\mathcal{B}$ has a finite subcover.
\end{thm}

\begin{proof}
Only the converse needs proof. Refine each member of an arbitrary open cover
by basic open subsets. A finite subcover from the resulting basic cover gives
a finite subcover from the original cover.
\end{proof}

\begin{thm}[Alexander's Subbase Theorem]
If $\mathcal{S}$ is a subbase for the topology of $X$, then $X$ is compact if
and only if every cover of $X$ by members of $\mathcal{S}$ has a finite
subcover.
\end{thm}

\begin{proof}
Only the reverse implication is nontrivial. Suppose every subbasic cover has
a finite subcover, but $X$ is not compact. By Zorn's lemma, there is a maximal
open cover $\mathcal{C}$ of $X$ having no finite subcover.

We claim that $\mathcal{C}\cap\mathcal{S}$ covers $X$. Otherwise choose
$x\in X$ that belongs to no subbasic member of $\mathcal{C}$. Since
$\mathcal{C}$ covers $X$, choose $U\in\mathcal{C}$ with $x\in U$. There are
$S_1,\ldots,S_n\in\mathcal{S}$ such that
\[
x\in S_1\cap\cdots\cap S_n\subseteq U.
\]
None of the $S_i$ belongs to $\mathcal{C}$. By maximality,
$\mathcal{C}\cup\{S_i\}$ has a finite subcover for each $i$; hence
$X=S_i\cup C_i$, where $C_i$ is a finite union of members of $\mathcal{C}$.
It follows that
\[
X=\bigcap_{i=1}^n(S_i\cup C_i)
\subseteq U\cup\bigcup_{i=1}^nC_i,
\]
which is a finite subcover drawn from $\mathcal{C}$, a contradiction. Thus
$\mathcal{C}\cap\mathcal{S}$ is a subbasic cover, so it has a finite
subcover, again a contradiction.
\end{proof}

\begin{de}
A topological space is connected if its only clopen subsets are $X$ and
$\varnothing$. A subset is connected if it is connected with its subspace
topology.
\end{de}

\begin{thm}
Let $X$ be a topological space.

(i) The continuous image of a connected set is connected.

(ii) If $\{E_i:i\in I\}$ is a family of connected subsets with
$E_i\cap E_j\neq\varnothing$ for all $i,j$, then $\bigcup_iE_i$ is connected.
If $(E_n)$ is a sequence of connected sets with
$E_n\cap E_{n+1}\neq\varnothing$, then $\bigcup_nE_n$ is connected.

(iii) The union of two intersecting connected sets is connected.

(iv) Every connected set is contained in a component; distinct components
are disjoint; and the components cover $X$.

(v) If $C$ is connected and
$C\subseteq Y\subseteq\operatorname{cl}C$, then $Y$ is connected.

(vi) The closure of a connected set is connected, and every component is
closed.
\end{thm}

\begin{proof}
The proofs are the same as in metric spaces; they use only inverse images of
open sets, relative topology, and the definition of closure.
\end{proof}

\begin{de}
A topological space $X$ is locally connected if, for each $x\in X$ and each
neighborhood $G$ of $x$, there is a connected open neighborhood $U$ of $x$
with $U\subseteq G$.
\end{de}

\begin{thm}
A topological space is locally connected if and only if it has a base
consisting of connected open sets.
\end{thm}

\begin{proof}
If $X$ is locally connected, let $\mathcal{B}$ be the collection of all
connected open subsets. Given $x\in G_1\cap G_2$ with $G_1,G_2\in\mathcal{B}$,
local connectedness applied to the neighborhood $G_1\cap G_2$ gives
$B\in\mathcal{B}$ with $x\in B\subseteq G_1\cap G_2$. Thus $\mathcal{B}$ is
a base. The converse follows immediately from the definition of a base.
\end{proof}

\begin{thm}
If $X$ is locally connected, then every component of $X$ is both open and
closed.
\end{thm}

\begin{proof}
Components are closed in every topological space. If $C$ is a component and
$x\in C$, local connectedness gives a connected open neighborhood $U$ of
$x$. Since $U\cap C\neq\varnothing$, the set $U\cup C$ is connected, so
maximality gives $U\subseteq C$. Therefore $C$ is open.
\end{proof}

\begin{thm}
If $X$ is locally connected and $f:X\to Z$ is continuous, open, and
surjective, then $Z$ is locally connected.
\end{thm}

\begin{proof}
Let $z\in Z$ and let $G$ be an open neighborhood of $z$. Choose $x\in X$
with $f(x)=z$. The set $f^{-1}(G)$ is an open neighborhood of $x$, so it
contains a connected open neighborhood $U$ of $x$. Then $f(U)$ is connected,
open, contains $z$, and is contained in $G$.
\end{proof}

\begin{eg}
(a) A connected space need not be locally connected; the topologist's sine
curve is a standard example.

(b) A locally connected space need not be connected; any discrete space with
more than one point is an example.
\end{eg}

\begin{ex}
If $X$ is compact, $Z$ is Hausdorff, and $f:X\to Z$ is a continuous
bijection, show that $f$ is a homeomorphism.
\end{ex}

\begin{proof}
If $F\subseteq X$ is closed, then $F$ is compact, so $f(F)$ is compact and
therefore closed in the Hausdorff space $Z$. Thus $f$ is a closed bijection,
and $f^{-1}$ is continuous.
\end{proof}

\begin{ex}
Use Alexander's Subbase Theorem to prove that a finite product of compact
spaces is compact.
\end{ex}

\begin{proof}
Use the standard subbase
\[
\{\pi_k^{-1}(U):1\leq k\leq n,\ U\text{ open in }X_k\}.
\]
Let $\mathcal{C}$ be a cover by such sets. For each $k$, let
\[
\mathcal{C}_k=\{U\subseteq X_k:\pi_k^{-1}(U)\in\mathcal{C}\}.
\]
At least one $\mathcal{C}_k$ covers $X_k$; otherwise choose
$x_k\notin\bigcup\mathcal{C}_k$ for every $k$, and the point
$(x_1,\ldots,x_n)$ would lie in no member of $\mathcal{C}$. Compactness of
that factor supplies a finite subcover, whose inverse images form a finite
subcover of the product.
\end{proof}

\begin{ex}
Give $\{0,1\}$ the discrete topology. Show that $X$ is connected if and only
if every continuous function $X\to\{0,1\}$ is constant.
\end{ex}

\begin{proof}
For a continuous $f:X\to\{0,1\}$, the sets $f^{-1}(\{0\})$ and
$f^{-1}(\{1\})$ are clopen and partition $X$. Thus a nonconstant $f$ gives a
separation, and every separation gives such a nonconstant function.
\end{proof}

\begin{ex}
Assume that $X_1,\ldots,X_n$ are nonempty. Show that the finite product
$X_1\times\cdots\times X_n$ is locally connected if and only if every
factor $X_k$ is locally connected.
\end{ex}

\begin{proof}
If the product is locally connected, each factor is locally connected because
its projection is continuous, open, and surjective. Conversely, let
$x=(x_1,\ldots,x_n)$ belong to a basic neighborhood
$G_1\times\cdots\times G_n$. Choose a connected open neighborhood
$U_k$ of $x_k$ contained in $G_k$. The finite product
$U_1\times\cdots\times U_n$ is connected and is an open neighborhood of $x$
contained in the given basic neighborhood.
\end{proof}

\subsection{Pathwise connectedness}

\begin{de}
A path in $X$ from $p$ to $q$ is a continuous function
$f:[0,1]\to X$ such that $f(0)=p$ and $f(1)=q$. A path with $p=q$ is a loop.
A space is path connected if every pair of points can be joined by a path.
\end{de}

\begin{thm}
If $f$ is a path from $a$ to $b$ and $g$ is a path from $b$ to $c$, then
\[
h(t)=
\begin{cases}
f(2t),&0\leq t\leq\tfrac12,\\
g(2t-1),&\tfrac12\leq t\leq1
\end{cases}
\]
is a path from $a$ to $c$.
\end{thm}

\begin{proof}
The two formulas agree at $t=1/2$, so the Pasting Lemma gives continuity.
\end{proof}

\begin{thm}
Let $X$ be a topological space.

(a) Every path-connected space is connected.

(b) A continuous image of a path-connected space is path connected; in
particular, this holds for every continuous surjection.

(c) If $\{E_i:i\in I\}$ is a collection of path-connected subsets with
$E_i\cap E_j\neq\varnothing$ for all $i,j$, then $\bigcup_iE_i$ is path
connected.

(d) If $(E_n)$ is a sequence of path-connected subsets with
$E_n\cap E_{n+1}\neq\varnothing$, then $\bigcup_nE_n$ is path connected.
\end{thm}

\begin{proof}
(a) If $X=A\cup B$ were a separation and a path joined a point of $A$ to a
point of $B$, its inverse images of $A$ and $B$ would separate $[0,1]$.

(b) Choose preimages of the two target points and compose a path between them
with the continuous map.

(c) Fix a common reference set $E_{i_0}$. Join each point to a point of
$E_{i_0}$ inside its own set, and concatenate with a path in $E_{i_0}$.

(d) Concatenate paths through the successive nonempty intersections; between
any two fixed points only finitely many sets are involved.
\end{proof}

\begin{thm}
Every path-connected subset is contained in a maximal path-connected subset,
called a path component. Distinct path components are disjoint, and the path
components cover the space.
\end{thm}

\begin{proof}
Use the union of all path-connected subsets containing a fixed point. All of
these sets meet at that point, so their union is path connected and maximal.
\end{proof}

\begin{de}
A space $X$ is locally path connected if, for each $x\in X$ and every
neighborhood $G$ of $x$, there is a path-connected open neighborhood $U$ of
$x$ with $U\subseteq G$.
\end{de}

Equivalently, $X$ has a base consisting of path-connected open sets.

\begin{thm}
If $X$ is locally path connected, then every open connected subset of $X$ is
path connected.
\end{thm}

\begin{proof}
Let $G\subseteq X$ be open. Every path component of $G$ is open in $G$:
through each of its points there is a path-connected open neighborhood
contained in $G$, and that neighborhood lies in the same path component. If
$G$ were connected and had more than one path component, one path component
and the union of the others would form a separation. Hence $G$ has one path
component.
\end{proof}

\begin{cor}
An open subset of $\mathbb{R}^q$ is connected if and only if it is path
connected.
\end{cor}

\begin{proof}
Euclidean space is locally path connected because open balls are path
connected.
\end{proof}

\textbf{Remark.} Path connectedness does not imply local path connectedness,
and local path connectedness does not imply path connectedness.

\subsection{Infinite products}

\begin{de}
Let $I$ be nonempty and let each $X_i$ be nonempty. Their Cartesian product is
\[
\prod_{i\in I}X_i
=\left\{x:I\to\bigcup_{i\in I}X_i:x(i)\in X_i\text{ for all }i\right\}.
\]
We write $x_i$ for $x(i)$. The coordinate projection is
$\pi_i(x)=x_i$.
\end{de}

\begin{thm}
Let $X$ be a set, let $X_i$ be topological spaces, and let
$f_i:X\to X_i$ be functions. The collection
\[
\mathcal{S}=\{f_i^{-1}(G):i\in I,\ G\text{ open in }X_i\}
\]
is a subbase for a topology $\mathcal{T}$ on $X$. This is the smallest
topology making every $f_i$ continuous. If the spaces $X_i$ are Hausdorff
and the functions separate points, then $\mathcal{T}$ is Hausdorff.
\end{thm}

\begin{proof}
The subbase construction gives a topology. Every $f_i$ is continuous by
definition. Any topology making every $f_i$ continuous must contain every
subbasic set and hence must contain $\mathcal{T}$. If $x\neq y$ and
$f_i(x)\neq f_i(y)$ for some $i$, disjoint neighborhoods of these two images
have disjoint inverse images separating $x$ and $y$.
\end{proof}

\begin{de}
The topology in the preceding theorem is the weak topology defined by the
family $\{f_i:i\in I\}$.
\end{de}

\begin{cor}
If $Z$ is a topological space and $g:Z\to X$, where $X$ has the weak topology
defined by $\{f_i\}$, then $g$ is continuous if and only if every
$f_i\circ g$ is continuous.
\end{cor}

\begin{proof}
One direction follows by composition. Conversely, the inverse image under
$g$ of a subbasic set $f_i^{-1}(G)$ is
$(f_i\circ g)^{-1}(G)$, which is open.
\end{proof}

\begin{de}
For a family of topological spaces $\{X_i:i\in I\}$, the weak topology on
$\prod_iX_i$ defined by the coordinate projections is the product topology.
The sets $\pi_i^{-1}(G)$ form the standard subbase.
\end{de}

\begin{thm}
Let $X=\prod_iX_i$ have the product topology.

(a) Every projection $\pi_i:X\to X_i$ is open.

(b) A map $g:Z\to X$ is continuous if and only if every coordinate map
$\pi_i\circ g$ is continuous.
\end{thm}

\begin{proof}
(a) The projection of a basic open set is open because all factors are
nonempty. Projections preserve unions.

(b) This is the weak-topology criterion.
\end{proof}

\begin{thm}
Let $(X_n,d_n)$ be metric spaces and put $X=\prod_{n=1}^\infty X_n$. Then the
product topology is induced by the metric
\[
D(x,y)=\sum_{n=1}^\infty2^{-n}
\frac{d_n(x_n,y_n)}{1+d_n(x_n,y_n)}.
\]
\end{thm}

\begin{proof}
Each summand is a bounded metric after composing $d_n$ with
$t\mapsto t/(1+t)$, so the convergent series defines a metric.

Let $G$ be a basic product neighborhood of $x$, restricting only the
coordinates $n_1,\ldots,n_k$, say
$d_{n_j}(x_{n_j},y_{n_j})<r_j$. Choose $\varepsilon>0$ so small that
\[
2^{n_j}\varepsilon<\frac{r_j}{1+r_j}
\qquad(1\leq j\leq k).
\]
Then $D(x,y)<\varepsilon$ implies all the required coordinate inequalities,
so a $D$-ball lies in $G$.

Conversely, let $B_D(x;\varepsilon)$ be a metric ball. Choose $N$ so that
$\sum_{n>N}2^{-n}<\varepsilon/2$. Choose $\delta>0$ such that
\[
\sum_{n=1}^N2^{-n}\frac{\delta}{1+\delta}<\frac{\varepsilon}{2}.
\]
Then the basic product neighborhood
\[
\bigcap_{n=1}^N\pi_n^{-1}(B_{d_n}(x_n;\delta))
\]
is contained in $B_D(x;\varepsilon)$. Thus the two topologies agree.
\end{proof}

\begin{thm}[Tychonoff's Theorem]
A product of nonempty topological spaces is compact if and only if every
factor is compact.
\end{thm}

\begin{proof}
If the product is compact, every factor is compact as the continuous image
under a surjective projection.

Conversely, suppose every $X_i$ is compact. By Alexander's Subbase Theorem it
is enough to consider a cover $\mathcal{C}$ by sets of the form
$\pi_i^{-1}(G)$. For each $i$, let
\[
\mathcal{C}_i=\{G\subseteq X_i:G\text{ is open and }
\pi_i^{-1}(G)\in\mathcal{C}\}.
\]
Some $\mathcal{C}_i$ must cover $X_i$. Otherwise choose
$x_i\in X_i\setminus\bigcup\mathcal{C}_i$ for every $i$; then
$x=(x_i)$ lies in no member of $\mathcal{C}$. Compactness of that factor gives
a finite subcover, whose inverse images form a finite subcover of the product.
\end{proof}

\begin{lem}
Let $X=\prod_iX_i$, fix $a\in X$, and set
\[
D=\{x\in X:x_i=a_i\text{ for all but finitely many }i\}.
\]
Then $D$ is dense in $X$.
\end{lem}

\begin{proof}
A nonempty basic open set restricts only finitely many coordinates. Choose a
point satisfying those restrictions and set every unrestricted coordinate
equal to the corresponding coordinate of $a$. The resulting point belongs to
both the basic open set and $D$.
\end{proof}

\begin{thm}
A product $\prod_iX_i$ is connected if and only if every factor $X_i$ is
connected.
\end{thm}

\begin{proof}
One direction follows from the projections. Conversely, fix $a\in X$. For
each finite $J\subseteq I$, let
\[
X_J=\{x\in X:x_i=a_i\text{ for }i\notin J\}.
\]
The set $X_J$ is homeomorphic to the finite product $\prod_{i\in J}X_i$,
which is connected. All the sets $X_J$ contain $a$, so their union $D$ is
connected. By the lemma, $D$ is dense; therefore
$X=\operatorname{cl}D$ is connected.
\end{proof}

\begin{thm}
A product $\prod_iX_i$ is path connected if and only if every factor is path
connected.
\end{thm}

\begin{proof}
One direction follows from the projections. Conversely, for $x,y\in X$,
choose a path $f_i:[0,1]\to X_i$ from $x_i$ to $y_i$ for every $i$. Define
$f(t)=(f_i(t))_{i\in I}$. Every coordinate $\pi_i\circ f=f_i$ is continuous,
so $f$ is continuous by the product criterion.
\end{proof}

\subsection{Nets}

\begin{de}
A directed set is a nonempty partially ordered set $I$ such that, for every
$i,j\in I$, there is a $k\in I$ with $i\leq k$ and $j\leq k$.

A net in a set $X$ is a function $x:I\to X$ from a directed set; it is usually
written $(x_i)_{i\in I}$.

A net in a topological space converges to $x$, written $x_i\to x$, if for
every neighborhood $G$ of $x$ there is an $i_0$ such that $x_i\in G$ whenever
$i\geq i_0$.

The net clusters at $x$, written $x_i\to_{\mathrm{cl}}x$, if for every
neighborhood $G$ of $x$ and every $j\in I$, there is an $i\geq j$ such that
$x_i\in G$.
\end{de}

\begin{thm}
(a) Every convergent net clusters at its limit.

(b) In a Hausdorff space, a net has at most one limit.
\end{thm}

\begin{proof}
Part (a) is immediate. For (b), suppose $x_i\to x$ and $x_i\to y$ with
$x\neq y$. Choose disjoint neighborhoods $G$ of $x$ and $H$ of $y$. The net
is eventually in each; choose an index larger than both corresponding
indices to obtain a contradiction.
\end{proof}

\begin{thm}
Let $X$ and $Z$ be topological spaces.

(a) A function $f:X\to Z$ is continuous at $x$ if and only if
$x_i\to x$ implies $f(x_i)\to f(x)$ for every net.

(b) If $f$ is continuous at $x$ and $x_i\to_{\mathrm{cl}}x$, then
$f(x_i)\to_{\mathrm{cl}}f(x)$.

(c) A subset $F\subseteq X$ is closed if and only if every convergent net of
points of $F$ has its limit in $F$.
\end{thm}

\begin{proof}
The forward implications follow directly from the neighborhood definition.
For the converse in (a), suppose $f$ is not continuous at $x$. There is a
neighborhood $W$ of $f(x)$ such that no neighborhood of $x$ is contained in
$f^{-1}(W)$. Direct the neighborhoods of $x$ by reverse inclusion, and for
each neighborhood $U$ choose
$x_U\in U\setminus f^{-1}(W)$. Then $x_U\to x$, but $f(x_U)$ is never in
$W$, a contradiction.

For (c), if $F$ is not closed and $x\in\operatorname{cl}F\setminus F$, direct
the neighborhoods of $x$ by reverse inclusion and choose $x_U\in U\cap F$.
Then $x_U\to x$.
\end{proof}

\begin{thm}
A topological space $X$ is compact if and only if every net in $X$ has a
cluster point.
\end{thm}

\begin{proof}
Suppose first that every net has a cluster point. If an open cover
$\mathcal{G}$ had no finite subcover, direct the finite subcollections
$\mathcal{F}\subseteq\mathcal{G}$ by inclusion and choose
\[
x_{\mathcal{F}}\in X\setminus\bigcup\mathcal{F}.
\]
Let $x$ be a cluster point and choose $G\in\mathcal{G}$ with $x\in G$. Taking
$\mathcal{F}_0=\{G\}$, the clustering property gives
$\mathcal{F}\supseteq\mathcal{F}_0$ with $x_{\mathcal{F}}\in G$, contradicting
the choice of $x_{\mathcal{F}}$.

Conversely, let $(x_i)$ be a net in compact $X$, and put
\[
F_j=\operatorname{cl}\{x_i:i\geq j\}.
\]
The family $(F_j)$ has the finite intersection property, so choose
$x\in\bigcap_jF_j$. Every neighborhood of $x$ meets every tail of the net;
hence $x$ is a cluster point.
\end{proof}

\begin{thm}
If $X$ is compact and a net $(x_i)$ has a unique cluster point $x$, then
$x_i\to x$.
\end{thm}

\begin{proof}
If the net did not converge to $x$, there would be a neighborhood $U$ of $x$
such that every tail contains a point in the closed set $X\setminus U$. The
closed sets
\[
(X\setminus U)\cap\operatorname{cl}\{x_j:j\geq i\}
\]
have the finite intersection property. Compactness would give a point
$y\in X\setminus U$ belonging to all of them, and $y$ would be another
cluster point.
\end{proof}

\begin{ex}
Let $\mathcal{S}$ be a subbase for $X$.

(a) Show that $x_i\to x$ if and only if, for every $S\in\mathcal{S}$ with
$x\in S$, the net is eventually in $S$.

(b) Find a subbase, a net, and a point $x$ such that the net is frequently in
every subbasic neighborhood of $x$ but does not cluster at $x$.
\end{ex}

\begin{proof}
(a) The forward implication is immediate. Conversely, a basic neighborhood
of $x$ is a finite intersection $S_1\cap\cdots\cap S_n$ of subbasic
neighborhoods. Choose a common upper bound for the finitely many indices after
which the net lies in each $S_k$.

(b) Let $X=\{1,2,3\}$ with the discrete topology and
\[
\mathcal{S}=\{\{1,2\},\{1,3\},\{2,3\}\}.
\]
This is a subbase. Let $x=1$ and let $x_n=2$ for even $n$ and $x_n=3$ for odd
$n$. The net is frequently in each subbasic neighborhood of $1$, but it never
enters the neighborhood $\{1\}$, so it does not cluster at $1$.
\end{proof}

\begin{ex}
Let $X=\prod_{\alpha\in A}X^\alpha$, let $(x_i)$ be a net in $X$, and write
$x_i=(x_i^\alpha)$ and $x=(x^\alpha)$.

(a) Show that $x_i\to x$ if and only if
$x_i^\alpha\to x^\alpha$ for every $\alpha$.

(b) Give a sequence in $\mathbb{R}^2$ whose two coordinate sequences cluster
at $x^1$ and $x^2$, respectively, but which does not cluster at
$(x^1,x^2)$.
\end{ex}

\begin{proof}
(a) This is the product continuity criterion applied to the coordinate
projections, or directly the subbase criterion for nets.

(b) Define
\[
x_n=
\begin{cases}
(1/n,0),&n\text{ even},\\
(1,1-1/n),&n\text{ odd}.
\end{cases}
\]
The first coordinate clusters at $0$, the second at $1$, but the sequence in
$\mathbb{R}^2$ does not cluster at $(0,1)$.
\end{proof}

\subsection{Quotient spaces}

\begin{de}
An equivalence relation $\sim$ on a set $X$ is reflexive, symmetric, and
transitive.
\end{de}

\begin{thm}
The equivalence classes of an equivalence relation form a partition of $X$.
Conversely, every partition of $X$ defines an equivalence relation by
declaring two points equivalent when they lie in the same member of the
partition.
\end{thm}

\begin{de}
The set of equivalence classes is denoted $X/\!\sim$. The natural map
$q:X\to X/\!\sim$ sends each point to its equivalence class and is called the
quotient map.
\end{de}

\begin{thm}
Let $X$ be a topological space with an equivalence relation $\sim$. The
collection
\[
\mathcal{U}=\{U\subseteq X/\!\sim:q^{-1}(U)\text{ is open in }X\}
\]
is a topology on $X/\!\sim$, called the quotient topology. The map $q$ is
continuous and surjective. The quotient need not be Hausdorff.
\end{thm}

\begin{proof}
Inverse images preserve the empty set, the whole set, arbitrary unions, and
finite intersections, so $\mathcal{U}$ is a topology. Continuity of $q$ is
built into its definition.
\end{proof}

\begin{thm}
Let $q:X\to X/\!\sim$ be a quotient map and let $f:X/\!\sim\to Z$. Then $f$
is continuous if and only if $f\circ q:X\to Z$ is continuous.
\end{thm}

\begin{proof}
One direction follows by composition. Conversely, if $U\subseteq Z$ is open,
then
\[
q^{-1}(f^{-1}(U))=(f\circ q)^{-1}(U)
\]
is open in $X$. By the definition of the quotient topology,
$f^{-1}(U)$ is open.
\end{proof}

\begin{ex}
Let $q:X\to X/\!\sim$ be an open quotient map. If $X$ is locally connected,
show that $X/\!\sim$ is locally connected.
\end{ex}

\begin{proof}
The map $q$ is continuous, open, and surjective, so the result follows from
the theorem on open continuous images of locally connected spaces.
\end{proof}

\begin{ex}
If $X$ is path connected and $\sim$ is any equivalence relation, show that
$X/\!\sim$ is path connected.
\end{ex}

\begin{proof}
The quotient map is continuous and surjective, and continuous images of
path-connected spaces are path connected.
\end{proof}

\begin{ex}
A semimetric (or pseudometric) on $X$ is a function
$\rho:X\times X\to[0,\infty)$ satisfying symmetry, the triangle inequality,
and $\rho(x,x)=0$, but distinct points may have distance zero. Give $X$ the
topology generated by the semimetric balls. Define $x\sim y$ when
$\rho(x,y)=0$.

(a) Show that $\sim$ is an equivalence relation.

(b) Show that
\[
d(q(x),q(y))=\rho(x,y)
\]
is a well-defined metric on $X/\!\sim$.

(c) Show that this metric induces the quotient topology.
\end{ex}

\begin{proof}
(a) Reflexivity and symmetry are immediate. If $\rho(x,y)=0$ and
$\rho(y,z)=0$, then the triangle inequality gives $\rho(x,z)=0$.

(b) If $x\sim x'$ and $y\sim y'$, then
\[
\rho(x,y)\leq\rho(x,x')+\rho(x',y')+\rho(y',y)=\rho(x',y'),
\]
and the reverse inequality follows symmetrically. Thus $d$ is well defined.
The metric axioms follow from those of $\rho$, and positive definiteness holds
because distinct equivalence classes have positive semidistance.

(c) For every $x\in X$ and $r>0$,
\[
q^{-1}(B_d(q(x);r))=B_\rho(x;r).
\]
Thus metric balls in the quotient are quotient-open. Conversely, if
$U\subseteq X/\!\sim$ is quotient-open and $q(x)\in U$, then
$q^{-1}(U)$ is an open neighborhood of $x$, so some
$B_\rho(x;r)$ lies in $q^{-1}(U)$. Hence
$B_d(q(x);r)\subseteq U$, proving that $U$ is metric-open.
\end{proof}
